\input{../tex-inputs/latex-leseansicht-vorspann}

\section[Arthur und Olga Schnitzler an Gustav Schwarzkopf, 9. 6. 1903]{L04430 Arthur und Olga Schnitzler an Gustav Schwarzkopf, 9. 6. 1903}
\nopagebreak\mylabel{L04430v}
\rehead{ }\normalsize\beginnumbering\briefempfaengerindex{Schwarzkopf, Gustav@\textsc{Schwarzkopf, Gustav}!zzzSchnitzler, Olga@\emph{von Olga Schnitzler}!1903-06-091@{9. 6. 1903}|(be}\briefempfaengerindex{Schwarzkopf, Gustav@\textsc{Schwarzkopf, Gustav}!zzzSchnitzler, Arthur@\emph{von Arthur Schnitzler}!1903-06-091@{9. 6. 1903}|(be}
\toendnotes[C]{\smallbreak\pagebreak[2]}
\correspDesc{Versand  durch Arthur Schnitzler,  am 9. 6. 1903 in Lugano
\newline{}Übermittlung  am 10. 6. 1903 in Lugano
\newline{}Zustellung  am 12. 6. 1903 in Wien 1/1 1
\newline{}Erhalt  durch Gustav Schwarzkopf am 12. 6. 1903 in Tiefer Graben 23}\toendnotes[C]{\smallbreak}
\Standort{DLA, A:Schnitzler, HS.1985.1.1897.}
\physDesc{Bildpostkarte, 159 Zeichen
\newline{}Handschrift Arthur Schnitzler: Bleistift, deutsche Kurrent
\newline{}Handschrift Olga Schnitzler: Bleistift, deutsche Kurrent
\newline{}Versand: 1) Stempel: »\nobreak{}\oindex{Lugano@\textbf{Lugano}|pwk}Lugano, 10 VI 03, 2\nobreak{}«.   2) Stempel: »\nobreak{}\oindex{I., Innere Stadt@\textbf{I., Innere Stadt}|pwk}Wien 1/1 1, 12. \textcolor{gray}{6}. \textcolor{gray}{03}, 8–9½V., Bestellt\nobreak{}«. }\toendnotes[C]{\smallbreak}\pstart{}{\pb}Herrn Guſtav Schwarzkopf\pend{}\pstart{}Wien I\oindex{I., Innere Stadt@\textbf{I., Innere Stadt}|pw}\pend{}\pstart{}Tiefer Graben 23\oindex{Wien@\textbf{Wien}!I., Innere Stadt@\textbf{I., Innere Stadt}!Tiefer Graben 23@\textbf{Tiefer Graben 23}|pw}.\pend{}{\bigskip}
\pstart
           \centering{}{\pb}\textcolor{gray}{\textbf{Melide\oindex{Melide@\textbf{Melide}|pw}, Bissone\oindex{Bissone@\textbf{Bissone}|pw} e Monte Generoso\oindex{Monte Generoso@\textbf{Monte Generoso}|pw}}}\pend
           \vspace{1em}
\pstart
           {\pb}Lugano\oindex{Lugano@\textbf{Lugano}|pw}{ }9. 6. 903.\pend
           \vspace{0.5em}
\pstart
           Herzliche Grüße!\pend
           \pstart Ihr \spacefill\mbox{A. S.}\pend{}
\pstart
           \noindent{}\label{T_L04430-1v}\edtext{Leider}{\lemma{\textnormal{\emph{Leider}}}\Cendnote{\textnormal{Daneben findet sich eine Wellenlinie mit Bleistift.
                     Eventuell ein Test des Schreibgeräts?}}}\label{T_L04430-1} regnets den dritten Tag.\pend
           \selectlanguage{ngerman}\vspace{1em}
\pstart
           \noindent{}{[}hs. O. Schnitzler:{]} Die beſten Grüße auch von mir!\pend
           \pstart \spacefill\mbox{OS.}\pend{}\selectlanguage{ngerman}\endnumbering\briefempfaengerindex{Schwarzkopf, Gustav@\textsc{Schwarzkopf, Gustav}!zzzSchnitzler, Olga@\emph{von Olga Schnitzler}!1903-06-091@{9. 6. 1903}|)be}\briefempfaengerindex{Schwarzkopf, Gustav@\textsc{Schwarzkopf, Gustav}!zzzSchnitzler, Arthur@\emph{von Arthur Schnitzler}!1903-06-091@{9. 6. 1903}|)be}\mylabel{L04430h}
\begin{anhang}
\end{anhang}\newcommand{\dateiname}{L04430}\newcommand{\titel}{Arthur und Olga Schnitzler an Gustav Schwarzkopf, 9. 6. 1903}\newcommand{\editorInnen}{Herausgegeben von Jahnke, SelmaMüller, Martin Anton}\input{../tex-inputs/latex-leseansicht-abspann}
